% Invented fixture for the checks fixed after the first real run of the skill (2026-11-24).
% Each pattern below is one that the scripts got wrong on a real manuscript. The physics is a placeholder,
% taken from the invented Mn3Sn example of examples/supplement/.
\documentclass[aip,reprint]{revtex4-2}
\begin{document}
\title{Placeholder title for the fixture}
\begin{abstract}
We simulate the current-induced reversal of tilted Mn$_3$Sn films and compare it with untilted films.
\end{abstract}
\maketitle

\section{Introduction}
Antiferromagnetic memory needs a reversal that a current alone can set~\cite{A}.
Earlier simulations of Mn$_3$Sn assumed an untilted kagome plane~\cite{B}.
These works do not evaluate how a tilt of the plane changes the threshold current.
Here we simulate tilted films and find that the threshold falls by a factor of two at a tilt of 30 degrees.
Sections~\ref{sec:model}--\ref{sec:comparison} give the model and the comparison, and Sec.~\ref{sec:conclusions} states the scope.

\section{Model}
\label{sec:model}
We use the energy
\begin{equation}
E=-K\cos^2\varphi-\mu_0 M H\cos\varphi,
\label{eq:first}
\end{equation}
and balance the torques on the octupole moment,
\begin{equation}
\tau_{\rm SOT}\sin\theta=\partial E/\partial\varphi.
\label{eq:second}
\end{equation}
Equation~\eqref{eq:second} defines the threshold, and Eq.~\eqref{eq:first} fixes the anisotropy that it must overcome.

\section{Comparison with untilted films}
\label{sec:comparison}
The tilt lowers the threshold by a factor of two at 30 degrees (Table~\ref{tab:values}).
Resolving the switching time in a film at room temperature also requires low damping.

\begin{table}[tb]
\caption{Threshold current density of the two films. Values in units of $10^{11}$ A/m$^2$.}
\label{tab:values}
\begin{ruledtabular}
\begin{tabular}{lc}
Film & $j_c$\\
tilted by 30 degrees & 6\\
tilted by 15 degrees & 12
\end{tabular}
\end{ruledtabular}
\end{table}

\section{Conclusions}
\label{sec:conclusions}
A tilt of the kagome plane lets a current alone reverse the octupole moment.
Our comparisons are between models.

\begin{acknowledgments}
Claude (Anthropic) assisted figure preparation and manuscript revision.
\end{acknowledgments}

\begin{thebibliography}{9}
\bibitem{A} A. Author, Placeholder Journal \textbf{1}, 1 (2020).
\bibitem{B} B. Author, Placeholder Journal \textbf{2}, 2 (2021).
\end{thebibliography}
\end{document}
